\chapter{Spatiotemporal mobility patterns in \textit{Bicing}} \label{ch:bicing_mobility_patterns}

This chapter summarizes the main mobility patterns observed in \textit{Bicing} during 2022, with a particular focus on the differences between \acp{m-b} and \acp{e-b}. This analysis provides context for the modeling work developed in Chapters~\ref{ch:results_mobility_and_battery_modeling} and~\ref{ch:results_failure_forecasting}, where mobility dynamics are linked to battery behavior and component wear.

\section{Methodology}

        \subsection{Data}

    The analysis is based on trip records provided by \textit{Bicing}, which encompass the following information:

    \begin{itemize}
            \item Starting and ending dates and times with second-level granularity
            \item Starting and ending stations
            \item Bike identifier
            \item Bike model (\ac{m-b} or \ac{e-b})
            \item Anonymized user identifier
    \end{itemize}

    The complete dataset covers trips from April 2019 to December 2022 (i.e., 3 years and 9 months), totaling 53 million trips (62\% \ac{m-b}, 38\% \ac{e-b}). However, this section uses only trips from 2022 to minimize the impact of the COVID-19 pandemic and to reflect the increasing use of \acp{e-b}, as detailed in Appendix~\ref{app_paper3:section_data_exploration}. Additionally, geographical data for the 519 stations, including latitude and longitude coordinates, are provided. 

    
    \subsection{Trips processing} \label{section_trips_processing}
  
    To study the mobility patterns and develop the \ac{pm} framework of Chapter~\ref{ch:results_failure_forecasting}, only trips with durations ranging from 2 to 60 minutes were taken into account, which constituted 99\% of all trips. Once filtered, various trip metrics were collected for the mobility analysis. While trip duration could be directly derived from the trips data, other measures required some additional steps. Obtaining the trip routes was not feasible since data only provided the starting and ending stations for each trip.

    As a consequence, the distances between stations were obtained using the OpenRouteService \ac{api} \cite{OpenRouteService_2023}, which can suggest recommended routes for bicycles based on \ac{osm} data~\cite{OpenStreetMap_2023}. These routes are generated by combining the suitability of streets for the chosen mode of transportation along with the fastest route option. 

    Then, using the trips' duration and distance their average speed was obtained. Furthermore, the cumulative trip inclinations were simplified by calculating the difference between the altitudes of the origin and destination stations. The specific station altitudes were obtained using the \ac{otd} \ac{api}~\cite{OpenTopoData_2023}.

    To determine whether the trip characteristics of \acp{m-b} and \acp{e-b} come from the same distribution, 5,000 samples from each subgroup were randomly selected. Initially, a Shapiro-Wilk test was applied to both samples to assess whether they followed a normal distribution. Subsequently, if both samples were found to be normally distributed, an Independent Sample T-Test was applied for comparison. In cases where one or both samples did not meet the normality assumption, the Kolmogorov-Smirnov test was used. All statistical tests were conducted with a significance level set at 0.01.

    To compare trip numbers of both bike models, several metrics were computed to avoid comparisons with absolute numbers. Initially, for each bike model, at all stations, the percentage of incoming (or outgoing) trips was determined by dividing the total trips for each model by the station's total trips. Then, two supplementary metrics were derived from these percentages: (1) the differences in the percentage of incoming (or outgoing) trips between \acp{e-b} and \acp{m-b}, and (2) the difference in the percentage of incoming and outgoing trips for each bike model.


\section{Results}

Both mechanical and electric transportation modes exhibited variability across time due to seasonal variations and holidays such as Easter week and summer holidays (Figure~\ref{fig:trips_characteristics}A). However, a noticeable differential trend emerged from summer onwards. Electric mobility experienced a rise in trip numbers while the mechanical one suffered a significant decrease. Despite some \acp{m-b} being upgraded to \acp{e-b} in the course of 2022, the increase in the number of \ac{e-b} rides couldn't be entirely attributed to this transformation, considering they comprised only 40\% of the bike fleet at their highest point. The \acp{e-b} experienced a rise in the mean daily trips per bike, whereas the \acp{m-b} saw a significant reduction. Consequently, the variations in the usage of both transportation modes can be attributed to the upgrade of \acp{m-b} into \acp{e-b} and the increasing preference of users for the electric mobility.

\begin{figure}[ht!]
    \centering
    \includegraphics[width=1\textwidth]{0_plots/paper3/main/panel_1.png}
    \caption[Trip patterns of \textit{Bicing} in 2022]{
        \textbf{Trip patterns of \textit{Bicing} in 2022.} 
        \textbf{(A)} Time series depicting the weekly trip counts alongside the average weekly number of trips per bike. 
        \textbf{(B)} Hourly time series illustrating the average trip counts throughout the week, with the area around the time series indicating the standard deviation. 
        \textbf{(C-F)} Distributions of trip characteristics, with vertical lines representing the mean of each subgroup.
    }
    \label{fig:trips_characteristics}
\end{figure}    

Trip counts present a strong weekly seasonality characterized by two distinct mainly daily patterns: weekdays and weekends (Figure~\ref{fig:trips_characteristics}B). On weekdays, three important trip count maxima are observed at 8:00, 14:00, and 18:00, with a gradual increase in trip numbers as the week progresses, peaking on Thursday. Fridays slightly deviate from the typical weekday pattern since they present similar trip count peaks at 14:00 and 18:00. In contrast, weekends are characterized by a 30\% mobility reduction, the absence of an 8:00 peak, and a shift in the 18:00 peak to 19:00. Notably, electric mobility consistently matches or surpasses the mechanical one across all days and hours despite its additional cost. Both transportation modes also differ in their trip characteristics (Figure~\ref{fig:trips_characteristics}C-F). Significant differences were found in the distributions of trip duration, distance, speed, and elevation (Section~\ref{section_trips_processing}), implying two distinct behaviors. Specifically, electric mobility is characterized by steeper inclines, longer durations, greater distances from the origin, and higher speeds compared to mechanical mobility.

Additionally, each mobility mode exhibits unique spatial mobility dynamics (Figure~\ref{fig:elevation_maps}). High-altitude stations predominantly receive \ac{e-b} trips, whereas low-altitude stations see the majority of incoming trips via \acp{m-b}. This reflects users' preference for electric mobility when doing physically demanding trips. While this pattern is clear for incoming trips, it becomes less pronounced for outgoing trips, since the proportional utilization of \acp{e-b} is not as dominant, even at high-altitude stations. This is attributed to the fact that the outgoing trips from a station do not indicate their destinations, resulting in a mix of trips to lower and higher stations. Nevertheless, given the overall preference for electric mobility among users, the inclination towards \acp{e-b} still remains evident. Furthermore, when analyzing the differences in the percentages of incoming and outgoing trips involving \acp{e-b} (Appendix~\ref{app_paper3:section_mobility_analysis}), it becomes evident that the prevalence of incoming \ac{e-b} trips is more pronounced at high-altitude stations. In contrast, at lower-altitude stations, the percentages of incoming trips made by \acp{e-b} become less substantial compared to the outgoing trips.

\begin{figure}[ht!]
    \centering
    \includegraphics[width=1\textwidth]{0_plots/paper3/main/panel_2.pdf}
    \caption[Trip flow percentages for e-bikes and m-bikes in 2022, segmented by trip elevation.]{
        \textbf{Trip flow percentages for \acp{e-b} and \acp{m-b} in 2022, segmented by trip elevation.} Each pair of maps displays the differences in \ac{e-b} and \ac{m-b} trip percentages for incoming and outgoing trips at each station. The left maps include all trips, whereas the remaining figures apply a filter based on trip elevation. Within each map, individual dots represent \ac{bss} stations. The dot size reflects the station's altitude, with larger dots indicating higher altitudes, and the dot color signifies differences in trip flow percentages. Grey dots denote stations that do not receive or send trips with the specified altitude filter.
    }
    \label{fig:elevation_maps}
\end{figure}

By segregating incoming and outgoing trips based on their elevation, this phenomenon becomes even more evident (Figure~\ref{fig:elevation_maps}). At stations that generate outgoing trips with an elevation gain of over 100 meters, electric mobility completely dominates the transportation mode. Consequently, the receiving stations at higher elevations predominantly receive \ac{e-b} trips. This prevalence of electric mobility reduces gradually as the elevation decreases until reaching those trips with elevations between -50 and 50 meters. In this range of elevations, mechanical mobility plays a much more prominent role, especially at low altitude stations, however, even at high-altitude stations, electric mobility remains relevant in both incoming and outgoing trips. Notably, even in scenarios with negative elevations exceeding -50 meters, electric mobility continues to be the preferred transport mode choice, albeit with reduced percentages. This phenomenon may be attributed to users' preference for an electric transportation mode on uphill rides, resulting in more \acp{e-b} accumulating at higher altitudes, and to potential elevation irregularities in paths between high-altitude stations.

Overall, these mode-specific temporal and topography-related mobility patterns provide the empirical context for the modeling frameworks developed in Chapters~\ref{ch:results_mobility_and_battery_modeling} and~\ref{ch:results_failure_forecasting}.
